\documentclass{ctexart}
\begin{document}

\begin{abstract}
本文针对连续介质中的传热过程建立机理模型，并用守恒关系、数值求解与误差检验形成闭环。
\end{abstract}
\keywords{机理模型；微分方程；数值离散；传热过程}

\section{问题重述}
题目要求在给定边界条件下计算温度场，并据此判断主要参数对终态的影响。

\section{问题分析}
热量在空间中的传递同时受到扩散系数和边界换热的制约，因此先写出守恒方程，再离散求解，并把结果回代到边界条件中检查。

\section{模型假设}
材料各向同性，研究时段内物性参数保持不变，外界温度按题目给定值处理。

\section{符号说明}
$T$ 表示温度，$k$ 表示导热系数，$t$ 表示时间。

\section{分问题模型建立与求解}
由能量守恒可得
\begin{equation}
  \frac{\partial T}{\partial t}=k\frac{\partial^2T}{\partial x^2}.
  \label{eq:heat}
\end{equation}
对式~\eqref{eq:heat} 采用有限差分离散，时间步长满足稳定性条件，离散思路参照经典传热文献\cite{heat}。表~\ref{tab:grid} 给出网格收敛结果。
\begin{table}[htbp]
  \centering
  \caption{网格收敛结果}\label{tab:grid}
  \begin{tabular}{cc}
    \hline
    网格数 & 终态误差 \\
    \hline
    100 & 0.018 \\
    200 & 0.009 \\
    \hline
  \end{tabular}
\end{table}

\section{结果分析}
计算表明，网格由 100 加密到 200 后，终态误差由 0.018 降至 0.009，变化方向与二阶空间离散的预期一致。由于边界附近的温度梯度最大，误差主要集中在换热边界，而内部节点已基本稳定。这说明当前网格能够支撑题目要求的数量级判断，同时也指出继续加密时应优先观察边界节点，而不是只报告全域平均值。

\section{模型检验}
将数值结果与一维稳态解析解对照，最大绝对误差为 0.009；把离散解回代到边界条件后，残差不超过 $1.1\times10^{-3}$，两项检验相互印证。

\section{灵敏度与稳健性分析}
在基准值上下扰动导热系数 5\%，终态平均温度的相对变化均小于 2\%。参数扰动未改变温度场的单调性，结论在该区间内保持稳健。

\section{模型评价与改进}
模型保留了守恒关系，参数含义清楚，且可用解析解复核；不足之处是未考虑温度对物性的反向影响。后续可引入温变参数并用自适应网格处理边界层。

\begin{thebibliography}{9}
\bibitem{heat} Carslaw H S, Jaeger J C. \emph{Conduction of Heat in Solids}. Oxford University Press, 1959.
\end{thebibliography}

\appendix
\section{程序结构}
程序依次读取参数、构造离散矩阵、推进时间步、计算残差并导出表格；每次运行保存网格数、步长和误差指标。

\end{document}
